\chapter*{Preface}
\addcontentsline{toc}{chapter}{Preface}

This book supports the module \textbf{Applied Physics in Industrial Design (APID)} through a single semester-long product development project. The activities are organised around hydroponics agriculture so that you can connect industrial design decisions to applied physics principles such as fluid flow, heat transfer, sensing, control, materials selection, and system integration.

\section*{Purpose}

The tutorial sequence is designed to help you move from an initial opportunity to a defensible product concept. Each session adds one part of the development process: identifying opportunities, understanding users, defining needs, translating needs into measurable specifications, generating concepts, selecting concepts, and communicating your final proposal.

\section*{Learning Approach}

The course uses project-based learning. You will work in teams, document your reasoning, and present intermediate outputs throughout the semester. This structure is intended to help you:
\begin{itemize}
    \item apply physics knowledge in a realistic design context;
    \item practise structured product development methods;
    \item build teamwork, documentation, and presentation skills; and
    \item reflect critically on technical and design decisions.
\end{itemize}

\section*{How to Use This Book}

Use each session as a working guide rather than as a reading-only text. Complete the pre-class tasks before the tutorial, bring your draft outputs to class for discussion, and revise the work after feedback. Unless stated otherwise, always confirm submission details, deadlines, and assessment updates through the university's \textbf{Institutional Teaching and E-Learning platform (ITEL)}.

\section*{Audience}

This material is written for Industrial Physics students who are developing confidence in both product development methods and professional communication. The book therefore emphasises clear templates, team workflows, and examples that show how engineering reasoning supports design decisions.

\section*{Acknowledgement}

This edition is still being refined. Feedback that improves clarity, consistency, or technical accuracy is valuable and should be shared with the module team through the usual course communication channels.
